\subsection{基于 NanoVDB 的 Hessian 差分提取}

在当前实现中，NanoVDB 默认采用 Trilinear 插值进行 SDF 查询\cite{museth2021nanovdb}。考虑到离散距离场在体素交界附近的高阶不连续性，本文并不直接对标量距离场 $\Phi$ 求 Hessian，而是先查询单位法向场
\begin{equation}
    \mathbf{n} = \frac{\nabla \Phi}{\|\nabla \Phi\|},
\end{equation}
再通过中心差分构造曲率张量
\begin{equation}
    \mathbf{W} \approx \nabla \mathbf{n}.
\end{equation}
实现中，差分步长固定取为
\begin{equation}
    h = 2 \cdot \text{voxel\_size}.
\end{equation}
这样得到的 $\mathbf{W}$ 对应于离散 SDF 法向流形上的 Weingarten 映射近似。

\subsection{张量坐标变换与 Patch 几何封装}

每个活动接触样本首先在模型坐标系下查询 $\Phi$、单位法向和曲率张量 $\mathbf{W}_M$。随后通过刚体姿态矩阵映射到世界坐标系：
\begin{equation}
    \mathbf{W}_W = \mathbf{R}_{WM}\mathbf{W}_M\mathbf{R}_{WM}^T.
\end{equation}

在活动集合构造阶段，系统首先依据预测法向相对速度执行 separating-velocity cutoff，滤除法向上已明显分离的候选样本；随后，其余候选样本再按空间半径阈值与法向夹角阈值聚类为局部 patch。patch 法向采用样本平均并重新归一化，而穿透深度 $\Phi$、梯度与 Hessian 等主几何量保留 deepest sample 的取值。经此处理后，每个 patch 被封装为可直接送入接触约束模块的局部几何代理。

\subsection{基于 OBB/BVH 的分层候选检测}

为降低连续滑移类算例中大规模从动体表面采样所带来的查询开销，当前实现还在样本侧引入了一个静态的 OBB/BVH broad-phase。该层次结构在从动体局部坐标系下构建，每个节点保存局部 OBB、节点代表采样点以及相应的包围半径；仿真过程中，仅需随刚体位姿将节点代表点映射到世界坐标，再映射到主动体 SDF 的模型坐标系下执行轻量 $\Phi$ 查询。

节点级剪枝遵循保守原则：实现首先利用主动体几何的包围球下界过滤明显不可能接触的节点，再结合节点代表点的 $\Phi$ 取值、节点包围半径以及一步运动余量进行 SDF 保守剪枝。只有未被剪除的叶节点样本才会进入后续的精确 $(\Phi,\nabla\Phi)$ 乃至 Hessian 查询。由此，OBB/BVH 仅改变候选枚举方式，而不改变叶节点内的接触样本构造、patch 聚类及曲率前馈公式，因此对最终接触结果保持无损。

该 broad-phase 默认服务于 \texttt{SlidingPatch} 路径中的连续滑移算例，如工业凸轮、偏心圆盘--滚子基准以及 Onset Stress Benchmark；对于 Twin Gear 这类已具备强时间相干局部扫描且接触切换更为紧凑的 \texttt{CompactDiscrete} 算例，当前默认仍保持关闭，仅作为离线可选优化。

\subsection{RHS 闭环注入与互补模块耦合}

对于每个活动 patch，系统根据当前步的接触点相对速度 $\mathbf{v}_{rel}$ 计算二阶曲率补偿
\begin{equation}
    C_{curve} = \frac{1}{2} (\Delta t)^2 (\mathbf{v}_{rel})^T \mathbf{W}_{W} \mathbf{v}_{rel},
\end{equation}
并构造修正后的有效接触间隙
\begin{equation}
    \Phi_{eff} = \Phi + C_{curve}.
\end{equation}

随后，该 $\Phi_{eff}$ 被直接写入接触约束容器，并作为 LCP/NSC 主方程右端项的一部分送入底层互补模块。由此，实现保持了对白盒寄生式注入结构的遵循，即不改写互补求解主链本身，而是通过高阶几何前馈修正接触约束输入。

算法 \ref{alg:slidingpatch} 与算法 \ref{alg:compactdiscrete} 分别概括了当前实现中两类局部接触表示的核心执行流程。二者共享同一底层 NSC/LCP 求解主链，但在候选样本筛选、patch 几何封装以及曲率项启用时机上存在本质区别。

\begin{algorithm}[htbp]
\caption{面向光滑连续接触的 \texttt{SlidingPatch} 流程}
\label{alg:slidingpatch}
\begin{algorithmic}[1]
\Require 从动体采样点 $\mathcal{S}$、上一时刻流形集合 $\mathcal{M}^{-}$、SDF 查询接口、模式 $\in \{\texttt{SdfFirstOrder},\texttt{SdfSecondOrder}\}$
\Ensure 带有有效间隙 $\Phi_{eff}$ 的活动接触 patch
\State 对 $\mathcal{S}$ 中每个样本查询 $(\Phi,\nabla \Phi)$
\State 剔除不满足 separating-velocity cutoff 的候选样本
\State 将剩余候选样本聚类为局部 patch
\State 将当前 patch 与 $\mathcal{M}^{-}$ 匹配，以保持 persistent manifold 身份
\State 通过 representative query 与局部几何拟合更新代表点
\State 对未通过 onset-aware filtering 的早激活候选进行拒绝
\If{当前模式为 \texttt{SdfSecondOrder}}
    \State 在代表几何处查询 Hessian / Weingarten 张量
    \State 计算 $C_{curve} \gets \frac{1}{2}(\Delta t)^2 \mathbf{v}_{rel}^{T}\mathbf{W}\mathbf{v}_{rel}$
    \State 置 $\Phi_{eff} \gets \Phi + C_{curve}$
\Else
    \State 置 $\Phi_{eff} \gets \Phi$
\EndIf
\State 将活动 patch 写入 NSC/LCP 接触容器
\end{algorithmic}
\end{algorithm}

\begin{algorithm}[htbp]
\caption{面向紧凑或奇异接触的 \texttt{CompactDiscrete} 流程}
\label{alg:compactdiscrete}
\begin{algorithmic}[1]
\Require 从动体采样点 $\mathcal{S}$、SDF 查询接口、模式 $\in \{\texttt{SdfFirstOrder},\texttt{SdfSecondOrder}\}$
\Ensure 保留主导几何的紧凑活动 patch 集合
\State 对 $\mathcal{S}$ 中每个样本查询 $(\Phi,\nabla \Phi)$
\State 剔除不满足 separating-velocity cutoff 的候选样本
\State 按空间半径阈值与法向夹角阈值对剩余候选进行聚类
\State 在每个 patch 内执行法向平均
\State 对 $\Phi$、梯度与 Hessian 保留主导/最深几何
\State 按预设的紧凑 patch 预算截断活动集合
\If{当前模式为 \texttt{SdfSecondOrder}}
    \State 在保留的 patch 几何上计算局部 $C_{curve}$
    \State 置 $\Phi_{eff} \gets \Phi + C_{curve}$
\Else
    \State 置 $\Phi_{eff} \gets \Phi$
\EndIf
\State 将紧凑活动 patch 写入 NSC/LCP 接触容器
\end{algorithmic}
\end{algorithm}
